\documentclass[11pt,a4paper]{article}
\usepackage[T1]{fontenc}
\usepackage[utf8]{inputenc}
\usepackage{lmodern,amsmath,amssymb,amsthm,mathtools,graphicx,booktabs,microtype}
\usepackage[margin=24mm,headheight=16pt]{geometry}
\usepackage{fancyhdr,hyperref,xcolor}
\definecolor{accent}{HTML}{91453A}
\hypersetup{colorlinks=true,urlcolor=accent,linkcolor=accent,pdftitle={Chapter 5: Identical Particles - Independent Worked Solutions}}
\pagestyle{fancy}\fancyhf{}\lhead{\small Quantum Mechanics / Chapter 5}\rhead{\small Kumar Azad}\cfoot{\thepage}
\setlength{\parindent}{0pt}\setlength{\parskip}{4pt}
\numberwithin{equation}{section}
\newcommand{\sh}{\hbar}
\newcommand{\dd}{\,\mathrm d}
\newcommand{\ii}{\mathrm i}
\newcommand{\ee}{\mathrm e}
\newcommand{\av}[1]{\langle #1\rangle}
\newcommand{\ket}[1]{\lvert#1\rangle}
\newcommand{\bra}[1]{\langle#1\rvert}
\newcommand{\bracket}[2]{\langle #1\mid#2\rangle}
\newcommand{\pnum}[2]{\subsection*{Problem #1 — #2}\addcontentsline{toc}{subsection}{Problem #1: #2}}
\newcommand{\result}[1]{\par\noindent\textcolor{accent}{\textbf{Result.}} #1\par}
\newcommand{\checkit}[1]{\par\noindent\textit{Interpretation / check.} #1\par}
\newcommand{\Aop}{\hat{\mathcal A}}
\newcommand{\Sop}{\hat{\mathcal S}}
\title{\textcolor{accent}{\large MATHEMATICAL PHYSICS · COURSEWORK COMPANION}\\[9pt]
\Huge\textbf{Identical Particles}\\[12pt]
\Large Quantum Mechanics · Chapter Five\\[5pt]
\normalsize Independent stepwise solutions to Problems 5.1–5.39}
\author{Kumar Azad\\\small\url{https://krazad0.github.io/QM-CH1-SOLUTION/chapter5/}}
\date{November 2026}
\begin{document}
\maketitle
\begin{abstract}
An independent graduate-style solutions companion for all thirty-nine numbered
exercises in Chapter 5 of Griffiths and Schroeter,
\emph{Introduction to Quantum Mechanics}, third edition.
The derivations develop reduced-mass separation, boson--fermion exchange,
Slater determinants, spin symmetries, helium and atomic configurations,
Fermi gases, magnetic response, periodic delta-function lattices,
and harmonic-chain phonons. Problem descriptions are original paraphrases:
consult the source textbook for precise wording and diagrams.
This is \textbf{not} an arXiv preprint, a peer-reviewed article,
or an official solution manual.
\end{abstract}
\tableofcontents
\section{General framework: exchange and symmetrization}
Let $P_{12}$ exchange every degree of freedom of two identical particles.
The spin--statistics rule (taken as an input in nonrelativistic quantum mechanics)
imposes
\begin{equation}
P_{12}\Psi=\begin{cases}+\Psi,&\text{identical bosons},\\
-\Psi,&\text{identical fermions}.\end{cases}
\end{equation}
For different orthonormal one-particle orbitals $a,b$,
\begin{equation}
\Psi_\pm(1,2)=\frac{\phi_a(1)\phi_b(2)\pm
\phi_b(1)\phi_a(2)}{\sqrt2}.
\end{equation}
For $N$ fermions, $\Psi_F=(N!)^{-1/2}\det[\phi_j(x_i)]$,
where $x_i$ includes position \emph{and} spin when appropriate.
A symmetric permanent with occupancy-dependent normalization describes
bosons. In three-dimensional free Fermi gases at zero temperature,
\begin{equation}
k_F=(3\pi^2n)^{1/3},\qquad
E_F=\frac{\sh^2k_F^2}{2m},\qquad
\frac{E_{\rm tot}}{N}=\frac35E_F
\end{equation}
for spin-$1/2$ degeneracy two.
We explicitly distinguish physical exchange symmetry from an actual force,
and use $E_F$ only within the zero-temperature, noninteracting model.

\section{Two-particle systems and exchange: Problems 5.1--5.8}
\pnum{5.1}{Two-particle separation into center-of-mass and relative motion}
Put $M=m_1+m_2$, $\mu=m_1m_2/M$,
$\mathbf r=\mathbf r_1-\mathbf r_2$ and
$\mathbf R=(m_1\mathbf r_1+m_2\mathbf r_2)/M$.
Solving these two linear vector equations gives
\begin{equation}
\boxed{\mathbf r_1=\mathbf R+\frac{\mu}{m_1}\mathbf r,\qquad
\mathbf r_2=\mathbf R-\frac{\mu}{m_2}\mathbf r}.
\end{equation}
The multivariable chain rule gives
$\nabla_1=(m_1/M)\nabla_R+\nabla_r
=(\mu/m_2)\nabla_R+\nabla_r$, and
$\nabla_2=(\mu/m_1)\nabla_R-\nabla_r$.
Substitute into $-\sh^2(\nabla_1^2/m_1+
\nabla_2^2/m_2)/2$. The mixed derivatives cancel:
\begin{equation}
\boxed{H=-\frac{\sh^2}{2M}\nabla_R^2
-\frac{\sh^2}{2\mu}\nabla_r^2+V(\mathbf r)}.
\end{equation}
For $\Psi(\mathbf R,\mathbf r)=F(\mathbf R)u(\mathbf r)$,
division by $Fu$ separates constants:
\begin{align}
-\frac{\sh^2}{2M}\nabla_R^2F&=E_{\rm CM}F,\\
\left[-\frac{\sh^2}{2\mu}\nabla_r^2+V(\mathbf r)\right]u&=E_{\rm rel}u,
\end{align}
where $\boxed{E=E_{\rm CM}+E_{\rm rel}}$.
This exactly maps a central two-body problem
onto a free center of mass and a single reduced-mass particle.

\pnum{5.2}{Reduced mass: isotope shifts, positronium and muonic hydrogen}
For an orbiting particle of mass $m$ and nuclear mass $M$,
$\mu=mM/(m+M)$. Coulomb energies scale with $\mu$,
so using $m$ in place of $\mu$ overestimates
the binding energy fractionally by
\begin{equation}
\frac{m-\mu}{\mu}=\frac mM.
\end{equation}
For ordinary hydrogen $m_e/m_p\simeq1/1836$,
giving $\boxed{0.054\%}$ (two significant figures).
For a Balmer transition $3\to2$,
\begin{equation}
\frac1\lambda=
R_\infty\frac{\mu}{m_e}
\left(\frac14-\frac19\right)
=\frac{5R_\infty\mu}{36m_e},
\quad
\lambda(M)=\frac{36m_e}{5R_\infty\mu(M)}.
\end{equation}
Because deuterium's nucleus is heavier,
its $\mu$ is larger, so its red line has a shorter
wavelength. Inserting $M_D\simeq2m_p$ gives
$\boxed{\lambda_H-\lambda_D\simeq0.18\ {\rm nm}}$.
For positronium, $\mu=m_e/2$, so
$\boxed{|E_1|=6.80\ {\rm eV}}$.
For muonic hydrogen $\mu_\mu=
m_\mu m_p/(m_\mu+m_p)$, with
$m_\mu=206.77m_e$, the reduced-mass ratio
$\mu_\mu/\mu_H\simeq186$;
the Lyman-$\alpha$ wavelength scales inversely
with $\mu$, giving
$\boxed{\lambda_\mu\simeq0.653\ {\rm nm}}$.
\checkit{The last value uses the nonrelativistic
point-nucleus spectrum; finite proton size and
QED corrections matter for precision muonic spectroscopy.}

\pnum{5.3}{Isotope splitting in the vibrational spectrum of HCl}
A diatomic oscillator has $\omega=\sqrt{k/\mu}$,
$\mu=m_Hm_{\rm Cl}/(m_H+m_{\rm Cl})$.
For approximate isotope masses $m_H=1$,
$m_{35}=35$, and $m_{37}=37$ (atomic units),
\begin{equation}
\mu_{35}=\frac{35}{36},\qquad
\mu_{37}=\frac{37}{38}.
\end{equation}
The frequencies are
$\nu_{35}/\nu_{37}=\sqrt{\mu_{37}/\mu_{35}}
=\sqrt{1332/1330}=1.0007516\ldots$,
so
\begin{equation}
\boxed{\Delta\nu\simeq7.51\times10^{-4}\nu}.
\end{equation}
The vibrational line for the heavier chlorine isotope
has the slightly lower frequency.

\pnum{5.4}{Normalization of bosonic and fermionic two-orbital states}
Take normalized orbitals $a,b$, with
$S=\bracket a b$. For
$\Psi_\pm=A[a(1)b(2)\pm b(1)a(2)]$,
\begin{equation}
\|\Psi_\pm\|^2=
|A|^2[2\pm2|S|^2].
\end{equation}
For orthogonal orbitals $S=0$,
$\boxed{A=1/\sqrt2}$ for both signs.
For the same normalized orbital $a=b$,
the symmetric bracket is $2a(1)a(2)$:
$\boxed{A=1/2}$.
The antisymmetric bracket is identically zero,
so no normalized fermion state exists.

\pnum{5.5}{Excited levels of two particles in a square well}
Let $\epsilon=\pi^2\sh^2/(2ma^2)$,
$\phi_n(x)=\sqrt{2/a}\sin(n\pi x/a)$.
For particles on $0<x_1,x_2<a$,
\begin{equation}
H=-\frac{\sh^2}{2m}
(\partial_{x_1}^2+\partial_{x_2}^2),
\quad E_{nm}=(n^2+m^2)\epsilon.
\end{equation}
For $n\ne m$, define
$\Phi^\pm_{nm}=
[\phi_n(1)\phi_m(2)\pm\phi_m(1)\phi_n(2)]/\sqrt2$.
Because the two product terms have the same sum
of one-body energies,
$H\Phi^-_{12}=5\epsilon\Phi^-_{12}$:
the fermionic ground state is verified.
Beyond the levels explicitly displayed in the
textbook example, the next two \emph{distinct}
levels are:
\begin{center}\begin{tabular}{lccc}\toprule
Statistics&Level&Wavefunction(s)&Degeneracy\\\midrule
Distinguishable&$8\epsilon$&$\phi_2(1)\phi_2(2)$&1\\
&$10\epsilon$&$\phi_1(1)\phi_3(2)$,
$\phi_3(1)\phi_1(2)$&2\\
Bosons&$8\epsilon$&$\phi_2(1)\phi_2(2)$&1\\
&$10\epsilon$&$\Phi^+_{13}$&1\\
Fermions&$10\epsilon$&$\Phi^-_{13}$&1\\
&$13\epsilon$&$\Phi^-_{23}$&1\\\bottomrule
\end{tabular}\end{center}
For fermions, $(2,2)$ is forbidden in this
spin-ignored model. Different ordered products
count as different states only when the
particles are distinguishable.

\pnum{5.6}{Exchange correction to mean squared separation in a well}
For $n\ne l$, the one-body moments are
\begin{align}
\av{x}_n&=a/2,&
\av{x^2}_n&=a^2\left(\frac13-\frac1{2\pi^2n^2}\right),\\
X_{nl}&=\int_0^a x\,\phi_n(x)\phi_l(x)\dd x\\
&=\frac{a}{\pi^2}\left[
\frac{(-1)^{n-l}-1}{(n-l)^2}
-\frac{(-1)^{n+l}-1}{(n+l)^2}\right].
\end{align}
Integrate $x\cos(k\pi x/a)$ by parts to
derive the last line; $X_{nl}=0$ for equal
parity, otherwise
$X_{nl}=-8anl/[\pi^2(n^2-l^2)^2]$.
For distinguishable particles,
$\av{(x_1-x_2)^2}_{D}=
\av{x^2}_n+\av{x^2}_l-2\av{x}_n\av{x}_l$:
\begin{equation}
\boxed{D_{nl}=a^2\left[
\frac16-\frac1{2\pi^2}
\left(\frac1{n^2}+\frac1{l^2}\right)\right]}.
\end{equation}
The exchange cross term in
$\av{x_1x_2}$ is $\pm|X_{nl}|^2$,
so
\begin{equation}
\boxed{\av{(x_1-x_2)^2}_{B}=D_{nl}-2|X_{nl}|^2,\quad
\av{(x_1-x_2)^2}_{F}=D_{nl}+2|X_{nl}|^2}.
\end{equation}
Bosons are closer on average in this
specific observable, and fermions farther apart.
This arises from exchange symmetry,
not an additional mechanical force.

\pnum{5.7}{Two oscillator orbitals: position densities and relative separation}
Let $b=\sqrt{\sh/(m\omega)}$,
$\phi_0(x)=\pi^{-1/4}b^{-1/2}
\exp[-x^2/(2b^2)]$,
and $\phi_1(x)=\sqrt2(x/b)\phi_0(x)$.
With one particle in each orbital,
\begin{equation}
\Psi_D=\phi_0(x_1)\phi_1(x_2),\qquad
\Psi_\pm=\frac{\phi_0(x_1)\phi_1(x_2)
\pm\phi_1(x_1)\phi_0(x_2)}{\sqrt2}.
\end{equation}
The distinguishable, boson, and fermion
probability surfaces are $|\Psi_D|^2$,
$|\Psi_+|^2$ and $|\Psi_-|^2$;
the last vanishes at $x_1=x_2$.
Since $\av{x^2}_0=b^2/2$,
$\av{x^2}_1=3b^2/2$ and
$\av{x}_{01}=b/\sqrt2$,
\begin{equation}
\boxed{\av{(x_1-x_2)^2}_D=2b^2,\quad
\av{(x_1-x_2)^2}_B=b^2,\quad
\av{(x_1-x_2)^2}_F=3b^2}.
\end{equation}
Introduce $R=(x_1+x_2)/2$ and
$r=x_1-x_2$, so
$x_1=R+r/2$, $x_2=R-r/2$.
The squared amplitudes factor into
a common Gaussian
$\exp[-2R^2/b^2-r^2/(2b^2)]$ times
$(R-r/2)^2$ for $D$, $4R^2$ for $B$,
or $r^2$ for $F$ (up to constants).
After integrating over $R$, put
$g_b(r)=\exp[-r^2/(2b^2)]/(\sqrt{2\pi}b)$.
The signed relative-displacement densities are
\begin{equation}
\boxed{p_D(r)=\tfrac12(1+r^2/b^2)g_b(r),\quad
p_B(r)=g_b(r),\quad
p_F(r)=(r^2/b^2)g_b(r)}.
\end{equation}
The distribution of distance $d=|r|\ge0$
is $\boxed{P_j(d)=2p_j(d)}$.
All three integrate to one, and the
second moments recover the results above.
Finally, the one-body number-density operator
$n(x)=\sum_{j=1}^2\delta(x-x_j)$ gives
the \emph{same} mean in every case:
\begin{equation}
\boxed{\av{n(x)}=
|\phi_0(x)|^2+|\phi_1(x)|^2}.
\end{equation}
Thus one-body density cannot distinguish
exchange symmetry here; two-body correlations can.

\pnum{5.8}{Three-particle bosonic permanent and fermionic determinant}
For three orthonormal orbitals $a,b,c$,
a particular distinguishable assignment is
$\boxed{\Psi_D=a(1)b(2)c(3)}$.
The completely symmetric state is
\begin{equation}
\boxed{\Psi_B=\frac1{\sqrt{3!}}
\sum_{\pi\in S_3}
a(\pi(1))b(\pi(2))c(\pi(3))}.
\end{equation}
Every one of the $3!$ terms is orthogonal
to all others, so the norm is $3!/3!=1$;
permutations merely reorder the sum.
The completely antisymmetric state is
\begin{equation}
\boxed{\Psi_F=\frac1{\sqrt{3!}}
\begin{vmatrix}
a(1)&b(1)&c(1)\\
a(2)&b(2)&c(2)\\
a(3)&b(3)&c(3)
\end{vmatrix}}.
\end{equation}
Exchanging two particle labels swaps
two rows and reverses the sign; repeated
orbitals give identical columns and zero.

\section{Spin, many-particle antisymmetry and atoms: Problems 5.9--5.16}
\pnum{5.9}{Including spin in low-lying two-particle well levels}
Take single-particle well energies $n^2\epsilon$,
$\epsilon=\pi^2\sh^2/(2ma^2)$, and
$\phi_n=\sqrt{2/a}\sin(n\pi x/a)$.
For fermions of spin $1/2$, antisymmetric
\emph{total} states pair a symmetric
spatial factor $\Phi^+_{nm}$ with the
spin singlet $\ket{0,0}$, and antisymmetric
$\Phi^-_{nm}$ with one of the three triplets
$\ket{1,M}$, $M=-1,0,1$.
When $n=m$, only the symmetric spatial
state survives, and there is just one
spin singlet. The first four distinct
energies and degeneracies are therefore
\begin{equation}
\boxed{
\begin{array}{c|c|c}
E/\epsilon&\text{states}&g\\\hline
2&\phi_1(1)\phi_1(2)\ket{0,0}&1\\
5&\Phi^+_{12}\ket{0,0}\text{ and }
\Phi^-_{12}\ket{1,M}&4\\
8&\phi_2(1)\phi_2(2)\ket{0,0}&1\\
10&\Phi^+_{13}\ket{0,0}\text{ and }
\Phi^-_{13}\ket{1,M}&4
\end{array}}.
\end{equation}
For hypothetical \emph{spin-1 bosons},
the two-spin exchange-symmetric sectors
have total $S=0$ (one state) and $S=2$
(five states), while the antisymmetric
sector has $S=1$ (three states).
A symmetric total bosonic wavefunction
therefore pairs $\Phi^+_{nm}$ with $S=0,2$
and $\Phi^-_{nm}$ with $S=1$.
For the same four energy levels, the
degeneracies are
\begin{equation}
\boxed{g(2\epsilon)=6,\quad
g(5\epsilon)=9,\quad
g(8\epsilon)=6,\quad
g(10\epsilon)=9}.
\end{equation}
These counts include spin and presume no
spin-dependent interaction.

\pnum{5.10}{Why three spin halves have no antisymmetric spin state}
A spin-$1/2$ particle has a two-dimensional
spin space $\mathbb C^2$. An entirely
antisymmetric spin state of $N$ particles
belongs to the exterior power
$\bigwedge^N\mathbb C^2$, of dimension
$\binom{2}{N}$. For $N=3$ that dimension
vanishes: $\boxed{\bigwedge^3\mathbb C^2=\{0\}}$.
Equivalently, every triple of arrows
$\uparrow,\downarrow$ contains two equal
arrows; swapping those must reverse the
sign of a completely antisymmetric state
but leaves the arrow assignment unchanged,
so its coefficient is zero.

For three spin-half fermions in a square
well, one spin orbital can be occupied
only once. The three lowest available
spin orbitals are $1\uparrow$, $1\downarrow$,
and either $2\uparrow$ or $2\downarrow$.
Let $\zeta_{n\sigma}(x,s)=\phi_n(x)\chi_\sigma(s)$.
The normalized ground states are the two
Slater determinants
\begin{equation}
\boxed{\Psi_{\sigma}=
\frac1{\sqrt{3!}}
\det\begin{pmatrix}
\zeta_{1\uparrow}(1)&\zeta_{1\downarrow}(1)&\zeta_{2\sigma}(1)\\
\zeta_{1\uparrow}(2)&\zeta_{1\downarrow}(2)&\zeta_{2\sigma}(2)\\
\zeta_{1\uparrow}(3)&\zeta_{1\downarrow}(3)&\zeta_{2\sigma}(3)
\end{pmatrix}},
\quad \sigma=\uparrow,\downarrow.
\end{equation}
Their energy is
$\boxed{E_{\rm ground}=2\epsilon+\epsilon+4\epsilon
=6\epsilon}$ and spin degeneracy is
$\boxed{2}$. In general the determinant
cannot be factored into one spatial
wavefunction times one spin wavefunction.

\pnum{5.11}{Laughlin wavefunction, Slater determinants, and correlation}
Ignoring an overall normalization, the
Laughlin state of $N$ electrons has the form
\begin{equation}
\Psi_q(z_1,\ldots,z_N)=
\left[\prod_{j<k}(z_j-z_k)^q\right]
\exp\left[-\frac12\sum_j|z_j|^2\right],
\qquad q=1,3,5,\ldots .
\end{equation}
Exchanging $z_j,z_k$ changes one factor's
sign by $(-1)^q=-1$; all other factors
are permuted in pairs. The Gaussian is
symmetric. Thus
$\boxed{P_{jk}\Psi_q=-\Psi_q}$.
For $q=1,N=3$, the Vandermonde formula is
\begin{equation}
\prod_{j<k}(z_j-z_k)=
-\det\begin{pmatrix}
1&z_1&z_1^2\\
1&z_2&z_2^2\\
1&z_3&z_3^2
\end{pmatrix}.
\end{equation}
Up to normalization and a global sign,
this is one Slater determinant formed from
single-particle orbitals
$\boxed{\varphi_m(z)\propto z^m\ee^{-|z|^2/2},
\quad m=0,1,2}$.
For $q=3,N=2$,
\begin{align}
(z_1-z_2)^3
&=(z_1^3-z_2^3)-3(z_1^2z_2-z_1z_2^2)\\
&=-\det\begin{pmatrix}1&z_1^3\\1&z_2^3\end{pmatrix}
-3\det\begin{pmatrix}z_1&z_1^2\\z_2&z_2^2\end{pmatrix}.
\end{align}
The state is a superposition of
determinants with different orbital
occupancies $(0,3)$ and $(1,2)$.
The Jastrow factor has a third-order
zero when particles meet, encoding
correlation not captured by a single
determinant. The normalized determinant
coefficients require the normalization
of the $z^m$ orbitals.

\pnum{5.12}{Symmetry projection for a general atomic eigenstate}
For $Z$ particles let $U_\pi$ permute their
position arguments. Since the identical-electron
Hamiltonian commutes with every permutation,
$[H,U_\pi]=0$. Define
\begin{equation}
\Sop=\frac1{Z!}\sum_{\pi\in S_Z}U_\pi,
\qquad
\Aop=\frac1{Z!}\sum_{\pi\in S_Z}
\operatorname{sgn}(\pi)\,U_\pi.
\end{equation}
If $H\psi=E\psi$, then
$H(\Sop\psi)=E(\Sop\psi)$ and
$H(\Aop\psi)=E(\Aop\psi)$.
These states must be normalized
if nonzero. If $\psi$ is symmetric
under exchange of two coordinates,
pair permutations that differ by this
exchange; contributions cancel,
giving $\boxed{\Aop\psi=0}$.
The one-electron spin Hilbert space
has only dimension two.
A completely antisymmetric spin state
of $Z$ spin-half electrons would lie
in $\bigwedge^Z\mathbb C^2$, which
vanishes whenever $Z>2$.
Thus $\boxed{Z>2:\text{ no nonzero
completely antisymmetric pure-spin factor}}$.
For atomic fermions the full position--spin
state is antisymmetric, but it typically
cannot be separated into pure spatial and
pure spin factors.

\pnum{5.13}{Autoionization energy and the hydrogenic helium-ion spectrum}
In the independent-electron approximation
of helium ($Z=2$), a hydrogenic electron
has $E_n=-54.4/n^2$ eV.
Two electrons both initially occupying
$n=2$ contribute
$E_i=2(-54.4/4)=-27.2$ eV.
If one falls to $1s$ while the other
is ejected, the final helium ion has
$E_{\mathrm{He}^+}=-54.4$ eV plus the
free electron's kinetic energy $K$.
Energy conservation
$-27.2=-54.4+K$ gives
\begin{equation}
\boxed{K=27.2\ {\rm eV}}.
\end{equation}
The helium ion $\mathrm{He}^+$ has one
electron and therefore the Coulomb
spectrum $E_n=-4(13.6\ {\rm eV})/n^2$.
Its radiative transitions obey
\begin{equation}
\boxed{\frac1\lambda=
4R_H\left(\frac1{n_f^2}-\frac1{n_i^2}\right),
\qquad n_i>n_f}.
\end{equation}
Finite nuclear-mass corrections change
the Rydberg constant slightly.
\checkit{The autoionization estimate ignores
electron--electron repulsion and requires
a physically allowed initial configuration;
it is a textbook idealization.}

\pnum{5.14}{Hypothetical exchange statistics and helium multiplets}
Physical electrons are identical fermions:
spatially symmetric orbitals pair with the
antisymmetric spin singlet, and spatially
antisymmetric orbitals pair with the symmetric
spin triplet. In ground-state helium both
electrons occupy $1s$, hence the spatial part
is symmetric and only a spin singlet is allowed.
If the electrons were instead hypothetical
identical \emph{bosons} with spin $1/2$,
the \emph{total} state would be symmetric.
The two-$1s$ ground configuration would
therefore have the symmetric spin triplet,
with $\boxed{3}$ spin projections.
For excited distinct orbitals $a,b$,
the symmetric spatial sector would pair with
the triplet, and the antisymmetric spatial
sector with the singlet. The former tends
to have larger direct Coulomb repulsion
than the latter; the relative para/ortho
pattern would be interchanged.
If the particles were \emph{distinguishable}
but retained the same masses and charges,
there would be no mandatory spin--space
symmetry pairing. Spatial symmetric and
antisymmetric eigenfunctions could each
occur with all four independent spin
configurations (one singlet and three triplets).
In the spin-independent Hamiltonian,
each spatial energy would have a fourfold
spin degeneracy, aside from any independent
spatial degeneracies. These hypotheticals
alter exchange rules, not the Coulomb operator.

\pnum{5.15}{Helium ground-state Coulomb repulsion integral}
In the uncorrelated $1s^2$ trial state
for charge $Z=2$,
\begin{equation}
\psi_0(\mathbf r_1,\mathbf r_2)
=\frac{Z^3}{\pi a^3}
\exp[-Z(r_1+r_2)/a],\qquad Z=2,
\end{equation}
where $a$ is the ordinary Bohr radius.
The normalized one-particle density is
$\rho(r)=Z^3\exp(-2Zr/a)/(\pi a^3)$.
The angular average of the Coulomb kernel
for fixed $r_1,r_2$ satisfies
\begin{equation}
\frac1{4\pi}\int
\frac{\dd\Omega_2}{|\mathbf r_1-\mathbf r_2|}
=\frac1{\max(r_1,r_2)};
\end{equation}
it follows by integrating
$[r_1^2+r_2^2-2r_1r_2u]^{-1/2}$
over $u\in[-1,1]$.
Set $b=2Z/a$.
The expectation is then
\begin{align}
I&=\int\dd^3r_1\dd^3r_2
\frac{\rho(r_1)\rho(r_2)}{|\mathbf r_1-\mathbf r_2|}\\
&=2b^6\int_0^\infty\dd r_1\,
r_1\ee^{-br_1}
\int_0^{r_1}\dd r_2\,r_2^2\ee^{-br_2}\\
&=\frac{5b}{16}=\frac{5Z}{8a}.
\end{align}
Thus for helium $\boxed{I=5/(4a)}$.
The interaction energy is
\begin{equation}
\Delta E=\frac{e^2}{4\pi\epsilon_0}I
=\frac54\frac{e^2}{4\pi\epsilon_0a}
\simeq\boxed{34.0\ {\rm eV}}.
\end{equation}
Adding to the independent-electron value
$E_0\simeq-108.8$ eV gives
$\boxed{E_{\rm trial}\simeq-74.8\ {\rm eV}}$,
compared with experimental
$E_{\rm He}\simeq-78.98$ eV.
This is a first-order / variational
estimate with an unoptimized $1s^2$ trial
state, not an exact helium solution.

\pnum{5.16}{Lithium: determinant ground state and degeneracy}
Neglect electron--electron interactions
for nuclear charge $Z=3$. The lowest
hydrogenic orbital has energy
$E_{1s}=-13.6Z^2=-122.4$ eV.
Two electrons occupy the distinct spin
orbitals $1s\uparrow$, $1s\downarrow$.
The third occupies $2s$ with
$E_{2s}=-122.4/4=-30.6$ eV.
Therefore
\begin{equation}
\boxed{E_{\rm Li}\simeq
2E_{1s}+E_{2s}=-275.4\ {\rm eV}}.
\end{equation}
For each $\sigma\in\{\uparrow,\downarrow\}$
the normalized wavefunction is
\begin{equation}
\boxed{
\Psi_{\sigma}(1,2,3)=
\frac1{\sqrt{3!}}\det
\begin{pmatrix}
1s\uparrow(1)&1s\downarrow(1)&2s\sigma(1)\\
1s\uparrow(2)&1s\downarrow(2)&2s\sigma(2)\\
1s\uparrow(3)&1s\downarrow(3)&2s\sigma(3)
\end{pmatrix}}.
\end{equation}
These are eigenstates of the sum of
hydrogenic one-body Hamiltonians and
are fully antisymmetric. The closed
$1s^2$ shell has zero total spin;
the $2s$ valence electron provides
a spin doublet:
$\boxed{g_{\rm ground}=2}$.

\section{Atomic configurations and free-electron solids: Problems 5.17--5.23}
\pnum{5.17}{Electronic configurations and atomic terms}
The Pauli principle allows two electrons of opposite
spin in each spatial orbital. Filling $1s$, $2s$,
and $2p$ subshells in the usual order gives:
\begin{center}\begin{tabular}{lll}\toprule
Element&Configuration&Ground term\\\midrule
H&$1s$&$\,{}^{2}S_{1/2}$\\
He&$1s^2$&$\,{}^{1}S_0$\\
Li&$1s^2\,2s$&$\,{}^{2}S_{1/2}$\\
Be&$1s^2\,2s^2$&$\,{}^{1}S_0$\\
B&$1s^2\,2s^2\,2p$&$\,{}^{2}P_{1/2}$\\
C&$1s^2\,2s^2\,2p^2$&$\,{}^{3}P_0$\\
N&$1s^2\,2s^2\,2p^3$&$\,{}^{4}S_{3/2}$\\
O&$1s^2\,2s^2\,2p^4$&$\,{}^{3}P_2$\\
F&$1s^2\,2s^2\,2p^5$&$\,{}^{2}P_{3/2}$\\
Ne&$1s^2\,2s^2\,2p^6$&$\,{}^{1}S_0$\\\bottomrule
\end{tabular}\end{center}
A term $\,{}^{2S+1}L_J$ specifies
spin $S$, orbital angular momentum $L$ and
total $J$. For H, He, Li, Be, the term values
follow immediately from their open or
closed $s$ shells. Boron's $2p$ electron
allows $\,{}^{2}P_{1/2}$ and $\,{}^{2}P_{3/2}$.
For carbon $2p^2$, antisymmetry allows
$\boxed{{}^{1}S_0,\ {}^{1}D_2,\ {}^{3}P_{0,1,2}}$.
For nitrogen $2p^3$, its allowed LS terms are
$\boxed{{}^{4}S_{3/2},\ {}^{2}D_{3/2,5/2},\
{}^{2}P_{1/2,3/2}}$.
Hund's energy-ordering rules select the
ground states.

\pnum{5.18}{Hund's three ordering rules}
(a) For excited $1s\,n\ell$ helium,
Hund's first rule favors a spin triplet
$S=1$ over a singlet $S=0$.
The required antisymmetric spatial
factor tends to reduce electron repulsion.
(b) Two $2p$ electrons in $L=2,M_L=2$
must both carry $m_\ell=1$; the orbital
top-of-ladder state is spatially symmetric.
It cannot accompany the symmetric spin
triplet, so at maximum $S=1$ carbon
has $L=1$, not $L=2$.
(c) Boron's $2p$ subshell is less than
half filled with $L=1,S=1/2$, hence
Hund's third rule selects
$\boxed{J=|L-S|=1/2}$,
the term $\,{}^{2}P_{1/2}$.
(d) For carbon, maximize $S=1$, use
allowed $L=1$, and choose $J=0$:
$\boxed{{}^{3}P_0}$.
For nitrogen $2p^3$, all three
$m_\ell=-1,0,+1$ orbitals are singly
occupied with parallel spin:
$S=3/2$, $M_L=0$, $L=0$, and
$J=3/2$, giving
$\boxed{{}^{4}S_{3/2}}$.
Hund rules are ordering principles
for LS-coupled many-electron atoms,
not exact solutions of the Coulomb problem.

\pnum{5.19}{Dysprosium's ground term}
A term $\,{}^{2S+1}L_J$ with label
$\,{}^{5}I_8$ has
$2S+1=5$, giving $\boxed{S=2}$.
The letter $I$ represents
$\boxed{L=6}$, and the subscript
gives $\boxed{J=8}$.
A likely dominant configuration is
$\boxed{[\mathrm{Xe}]\,4f^{10}\,6s^2}$.
The closed $6s^2$ pair contributes
zero net $L$ and $S$.
The open $4f^{10}$ shell supports
$S=2,L=6$, and being more than half
filled favors $J=L+S=8$.

\pnum{5.20}{Mean energy in a filled three-dimensional Fermi sphere}
At zero temperature, spin-$1/2$
electrons fill $k\le k_F$. Counting
two spin orientations,
\begin{equation}
N_e=2\frac{V}{(2\pi)^3}4\pi
\int_0^{k_F}k^2\dd k
=\frac{Vk_F^3}{3\pi^2}.
\end{equation}
Their total energy is
\begin{equation}
E_{\rm tot}=2\frac{V}{(2\pi)^3}4\pi
\int_0^{k_F}
\frac{\sh^2k^2}{2m}k^2\dd k
=\frac{V\sh^2k_F^5}{10\pi^2m}.
\end{equation}
Dividing gives
\begin{equation}
\boxed{\frac{E_{\rm tot}}{N_e}
=\frac35\frac{\sh^2k_F^2}{2m}
=\frac35E_F}.
\end{equation}

\pnum{5.21}{Copper: Fermi energy, speed, temperature and pressure}
The copper mass density is
$\rho_m=8.96\times10^3\,{\rm kg/m^3}$,
molar mass $0.0635\,{\rm kg/mol}$,
with one free electron per atom.
The conduction-electron density is
\begin{equation}
n=\frac{\rho_m N_A}{0.0635\,{\rm kg/mol}}
\simeq8.50\times10^{28}\ {\rm m}^{-3}.
\end{equation}
Using
$E_F=\sh^2(3\pi^2n)^{2/3}/(2m_e)$,
\begin{equation}
\boxed{E_F\simeq7.03\ {\rm eV}}.
\end{equation}
The Fermi speed is
$\boxed{v_F=\sqrt{2E_F/m_e}\simeq
1.57\times10^6\,{\rm m/s}\simeq0.0052c}$,
so the nonrelativistic assumption is safe.
The Fermi temperature and degeneracy pressure are
\begin{equation}
\boxed{T_F=E_F/k_B\simeq8.15\times10^4\ {\rm K}},
\quad
\boxed{P=\frac25nE_F\simeq3.83\times10^{10}\ {\rm Pa}}.
\end{equation}
$E_F$ must be converted to Joules
in the pressure calculation.
These are noninteracting zero-temperature
electron-gas estimates.

\pnum{5.22}{Helium-3 Fermi temperature}
A helium-3 atom has spin one-half
and mass $m_3\simeq5.01\times10^{-27}$ kg.
For density $82\,{\rm kg/m^3}$,
the particle density is
$n=82/m_3\simeq1.64\times10^{28}\,{\rm m^{-3}}$.
For spin degeneracy two,
\begin{equation}
\boxed{T_F=
\frac{\sh^2}{2m_3k_B}(3\pi^2n)^{2/3}
\simeq3.1\ {\rm K}}.
\end{equation}
This ideal-gas estimate omits
liquid helium-3 interactions.

\pnum{5.23}{Bulk modulus of the degenerate electron gas}
At fixed electron count $N_e$,
$n=N_e/V$ and
$P=Cn^{5/3}=CN_e^{5/3}V^{-5/3}$.
Therefore
\begin{equation}
\frac{\dd P}{\dd V}=-\frac53\frac{P}{V},
\qquad
\boxed{B=-V\frac{\dd P}{\dd V}
=\frac53P}.
\end{equation}
Using the copper pressure in 5.21,
\begin{equation}
\boxed{B_{\rm Cu}^{\rm ideal}
\simeq6.38\times10^{10}\ {\rm Pa}}.
\end{equation}
The textbook's comparison value is
$1.34\times10^{11}\,{\rm Pa}$;
the discrepancy is expected from
interactions and band structure
neglected by the gas model.

\section{Periodic potentials and band structure: Problems 5.24--5.28}
\pnum{5.24}{Wavefunctions in a periodic delta-function comb}
Inside a cell $0<x<a$ the general solution
is $\psi(x)=A\sin(kx)+B\cos(kx)$.
Continuity and the Bloch condition imply
$A\sin(ka)=[\ee^{\ii qa}-\cos(ka)]B$.
For $\sin(ka)\ne0$, substitute and
use $\sin(ka)\cos(kx)-\cos(ka)\sin(kx)
=\sin[k(a-x)]$:
\begin{equation}
\boxed{\psi(x)=C\{
\sin(kx)+\ee^{-\ii qa}\sin[k(a-x)]\}},
\qquad 0\le x\le a,
\end{equation}
where $C=A/[\ldots]$ absorbs an
unimportant constant and phase.
Equivalently this follows by choosing
$C=B\ee^{\ii qa}/\sin(ka)$.
At a band boundary $z=ka=j\pi$,
the ratio used in this derivation is
singular. The correct state has zeros
at every delta spike:
\begin{equation}
\boxed{\psi(x)\propto
\sin(j\pi x/a)}
\end{equation}
extended globally as a standing wave.
Its derivative is continuous across
each spike because the jump
$2m\alpha\psi/\sh^2$ vanishes there.
The Bloch phase at this boundary is
$\ee^{\ii qa}=(-1)^j$.

\pnum{5.25}{Lower edge of the first repulsive delta-comb band}
Define $z=ka$ and
$\beta=m\alpha a/\sh^2$.
The Bloch dispersion equation is
\begin{equation}
\boxed{\cos(qa)=f(z)=
\cos z+\beta\frac{\sin z}{z}}.
\end{equation}
For $\beta=10$, $f(0)=11$,
which is forbidden because $|\cos qa|\le1$.
The lowest band starts at the first
positive root of
$f(z)=1$, which can be found with
bisection or Brent's method.
Numerically, $\boxed{z_0=2.627675433\ldots}$.
The energy is $E=\sh^2z^2/(2ma^2)$.
Since $\alpha/a=1$ eV and
$\beta=m\alpha a/\sh^2=10$,
$\sh^2/(2ma^2)=\alpha/(2\beta a)=0.05$ eV.
Therefore
\begin{equation}
\boxed{E_{\rm band,bottom}=
0.05z_0^2\ {\rm eV}\simeq0.345\ {\rm eV}}
\end{equation}
(three significant digits).

\pnum{5.26}{Attractive delta comb and negative-energy Bloch bands}
For an attractive lattice set
$\beta=-|\beta|$ in the same
positive-energy dispersion function:
\begin{equation}
f(z)=\cos z-|\beta|\frac{\sin z}{z},
\qquad z=ka>0.
\end{equation}
For negative energies, let
$\kappa=\sqrt{-2mE}/\sh>0$
and $t=\kappa a$. Analytic continuation
$k=\ii\kappa$ gives
\begin{equation}
\boxed{\cos(qa)=
\cosh t-|\beta|\frac{\sinh t}{t}}.
\end{equation}
Allowed states have the right-hand
side in $[-1,1]$. For $|\beta|=1.5$,
$f(E=0)=1-1.5=-0.5$, so the
lowest band crosses $E=0$, containing
both negative- and positive-energy
states. As $E\to-\infty$,
$\cosh t-|\beta|\sinh t/t\to\infty$
and negative energies become forbidden.
There is therefore a finite lower edge
for the bound portion of the band.
A periodic chain of $N$ cells has
$\boxed{N\text{ spatial Bloch states
per allowed band}}$ with periodic
boundary conditions. Only some may
have $E<0$ in this parameter regime;
the rest lie in its positive-energy
continuation.

\pnum{5.27}{Bloch momentum reversal and energy degeneracies}
For a ring of $N$ identical cells,
$q_n=2\pi n/(Na)$, with
$n=0,\ldots,N-1$.
The dispersion relation depends
on $q$ through $\cos(qa)$.
Hence
\begin{equation}
E(q)=E(-q)
\end{equation}
within any fixed band.
The modes $n$ and $N-n$ have the
same cosine and normally form a pair
of degenerate but independent traveling
Bloch states. Exceptions occur when
the two modes coincide modulo the
reciprocal lattice vector:
\begin{equation}
q\equiv-q\pmod{2\pi/a}.
\end{equation}
This means $2n\equiv0\pmod N$.
Thus
\begin{equation}
\boxed{n=0\text{ always};\qquad
n=N/2\text{ also if }N\text{ is even}}.
\end{equation}
Those time-reversal invariant momenta
are not forced into distinct $\pm q$
pairs. At zero and zone-edge
momenta there can be nondegenerate
band-edge eigenfunctions.

\pnum{5.28}{Reproducible dispersion plot of the delta comb}
With $m=\sh=a=\alpha=1$,
\begin{equation}
E=\frac{k^2}{2},\qquad
\cos q=\cos k+\frac{\sin k}{k}\equiv f(k).
\end{equation}
For each sampled $E\in(0,30]$,
calculate $k=\sqrt{2E}$ and $f(k)$.
If $|f(k)|\le1$, define
$q=\arccos f(k)$ and add the two
states $(q,E)$ and $(-q,E)$.
Otherwise the energy belongs to a
forbidden gap and should not be
plotted as a Bloch eigenvalue.
For the interval $[-\pi,\pi]$ this
constructs the bands in the first
Brillouin zone. The script
\texttt{chapter5/scripts/verify.py}
produces \texttt{figures/bands.pdf}
using $30{,}000$ energy samples.
\checkit{The plotted curve is an implicit
dispersion relation, not the free-electron
parabola; gaps occur where $|f(k)|>1$.}

\section{Quantum statistics and magnetic response: Problems 5.29--5.34}
\pnum{5.29}{Three particles distributed among three orbitals}
There are three distinguishable particles and
three available orthonormal one-body states.
Each labeled particle independently has
three choices, so
$\boxed{\mathcal N_D=3^3=27}$.
For three identical bosons, an arrangement
is specified by nonnegative occupancies
$n_a+n_b+n_c=3$. Stars-and-bars gives
\begin{equation}
\boxed{\mathcal N_B=
\binom{3+3-1}{3}=\binom53=10}.
\end{equation}
For three identical fermions with no additional
spin degree of freedom, the Pauli principle
allows at most one particle in each orbital,
so all three must be singly occupied:
\begin{equation}
\boxed{\mathcal N_F=\binom33=1}.
\end{equation}
The last state is the unique antisymmetric
Slater determinant up to phase.

\pnum{5.30}{Fermi energy in a two-dimensional square well}
For a large two-dimensional well, count
states in momentum space in the
thermodynamic limit. With area $A$,
each $k$-state occupies $(2\pi)^2/A$
and each accepts two spin orientations.
A filled disk of radius $k_F$ has
\begin{equation}
N_e=2\frac{A}{(2\pi)^2}\pi k_F^2
=\frac{Ak_F^2}{2\pi}.
\end{equation}
Let $\sigma=N_e/A$ be the areal density.
Then
$\boxed{k_F=\sqrt{2\pi\sigma}}$,
and
\begin{equation}
\boxed{E_F=\frac{\sh^2k_F^2}{2m}
=\frac{\pi\sh^2\sigma}{m}}.
\end{equation}
Boundary effects vanish relative to
the area in the macroscopic limit.

\pnum{5.31}{Cohesive-energy estimate including spin in three dimensions}
In the textbook's idealized box analogy,
take $N$ isolated monovalent atoms, each
hosting an electron in a cubic box of
side $a$. The individual ground-state
energy is
\begin{equation}
E_{\rm atom}=
\frac{\pi^2\sh^2}{2ma^2}(1^2+1^2+1^2)
=\frac{3\pi^2\sh^2}{2ma^2}.
\end{equation}
The separated total is $NE_{\rm atom}$.
After combining into a cubic metal
of volume $V=Na^3$, the same $N$
electrons occupy free-electron states
with twofold spin degeneracy.
Their number density is
$n=N/V=1/a^3$, giving
\begin{equation}
E_F=\frac{\sh^2}{2ma^2}(3\pi^2)^{2/3},
\quad
E_{\rm metal}=\frac35 NE_F.
\end{equation}
Define the cohesive-energy gain
as the separated energy minus the
metal energy. Per atom,
\begin{equation}
\boxed{\frac{\Delta E_{\rm coh}}N=
\frac{\sh^2}{2ma^2}\left[
3\pi^2-\frac35(3\pi^2)^{2/3}\right]}.
\end{equation}
It is positive in this model, indicating
an energetic advantage to delocalizing
electrons across a much larger volume.
The comparison intentionally ignores
the detailed ionic binding potentials
and is only an illustrative toy model.

\pnum{5.32}{Fermi gas with unequal spin populations}
Let $N_\pm$ be the spin-labeled counts
in volume $V$ and $n_\pm=N_\pm/V$.
Each spin species has degeneracy one,
so its Fermi wavevector is
$k_{F,\pm}=(6\pi^2n_\pm)^{1/3}$.
Integrating the kinetic energy for
each species gives
\begin{equation}
\boxed{E_{\rm tot}=
\frac{3\sh^2}{10m}(6\pi^2)^{2/3}
V(n_+^{5/3}+n_-^{5/3})}.
\end{equation}
For equal populations $n_\pm=\rho/2$,
the result reduces to
$E_{\rm tot}/V=
(3/10)(\sh^2/m)(3\pi^2)^{2/3}
\rho^{5/3}$, the unpolarized formula.
With magnetization
$M=-\mu_B(n_+-n_-)$,
define $\eta=M/(\rho\mu_B)$, so
$n_\pm=\rho(1\mp\eta)/2$.
Substitute:
\begin{equation}
\frac{E_{\rm tot}}V=
\frac{3}{10}\frac{\sh^2}{m}
(3\pi^2)^{2/3}\rho^{5/3}
\frac{(1+\eta)^{5/3}+
(1-\eta)^{5/3}}2.
\end{equation}
The binomial expansion yields
$\frac12[(1+\eta)^{5/3}+
(1-\eta)^{5/3}]=
1+\frac59\eta^2+O(\eta^4)$.
Thus
\begin{equation}
\boxed{\frac{E_{\rm tot}}V=
\frac{\sh^2(3\pi^2\rho)^{5/3}}
{10\pi^2m}
\left[1+\frac59
\left(\frac{M}{\rho\mu_B}\right)^2
+O(M^4)\right]}.
\end{equation}
The positive quadratic term makes
$M=0$ the kinetic-energy minimum.

\pnum{5.33}{Pauli paramagnetism and aluminum susceptibility}
For a weak uniform field along $z$,
the Zeeman energy density is $-MB$.
From Problem 5.32, the kinetic
energy density for small $M$ is
\begin{equation}
\frac{E_{\rm kin}}V
=\frac35\rho E_F+
\frac{E_F}{3\rho\mu_B^2}M^2+O(M^4).
\end{equation}
Minimize
$\mathcal E(M)=E_{\rm kin}/V-MB$:
\begin{equation}
\frac{\partial\mathcal E}{\partial M}
=\frac{2E_F}{3\rho\mu_B^2}M-B=0,
\qquad
\boxed{M=\frac{3\rho\mu_B^2}{2E_F}B}.
\end{equation}
With susceptibility defined by
$\chi=\mu_0\dd M/\dd B$,
\begin{equation}
\boxed{\chi_P=
\frac{3\mu_0\rho\mu_B^2}{2E_F}},
\quad
E_F=\frac{\sh^2}{2m_e}(3\pi^2\rho)^{2/3}.
\end{equation}
For aluminum $\rho=18.1\times10^{22}
\,{\rm cm^{-3}}=1.81\times10^{29}\,{\rm m^{-3}}$,
substitution gives
$\boxed{\chi_P\simeq1.2\times10^{-5}}$.
The textbook's comparison value is
$\chi_{\rm exp}=22\times10^{-6}$.
This is spin-only Pauli susceptibility;
orbital diamagnetism and interactions
are intentionally omitted.

\pnum{5.34}{Stoner criterion in a pair-counting interaction model}
The model assigns extra energy $U$ to
each opposite-spin pair:
\begin{equation}
\Delta E=UN_+N_-.
\end{equation}
Since $N_++N_-=\rho V$ and
$N_+-N_-=-MV/\mu_B$,
\begin{equation}
\boxed{\Delta E=
\frac{UV^2}{4}
\left(\rho^2-\frac{M^2}{\mu_B^2}\right)}.
\end{equation}
From Problem 5.32 the kinetic
contribution quadratic in $M$ is
$V E_F M^2/(3\rho\mu_B^2)$.
Thus the total quadratic coefficient is
\begin{equation}
\frac{V}{\mu_B^2}
\left(\frac{E_F}{3\rho}
-\frac{UV}{4}\right).
\end{equation}
The unpolarized state ceases to be a local
minimum when this becomes negative, giving
\begin{equation}
\boxed{U>U_c=\frac{4E_F}{3\rho V}
=\frac{4E_F}{3(N_++N_-)}}.
\end{equation}
\checkit{$U$ is the energy per
opposite-spin \emph{pair}. For an
extensive thermodynamic interaction,
one normally scales $U\propto1/V$.
This simple instability condition
is the stated textbook model,
not a complete theory of ferromagnetism.}

\section{Degeneracy pressure and phonons: Problems 5.35--5.39}
\pnum{5.35}{Nonrelativistic white-dwarf degeneracy pressure}
Let the star contain $N$ nucleons of mass
$M$, and $N_e=dN$ electrons of mass $m$.
Assume a uniformly dense sphere with
$V=4\pi R^3/3$. In a nonrelativistic
spin-$1/2$ Fermi gas,
\begin{align}
E_e&=\frac{3\sh^2}{10m}
(3\pi^2)^{2/3}(dN)^{5/3}V^{-2/3}\\
&=\frac{3}{10}
\left(\frac{9\pi}{4}\right)^{2/3}
\frac{\sh^2}{m}
\frac{d^{5/3}N^{5/3}}{R^2}.
\end{align}
A uniform Newtonian sphere of mass
$NM$ has gravitational potential energy
\begin{equation}
\boxed{E_g=-\frac35\frac{G(NM)^2}{R}}.
\end{equation}
The total energy has the form
$E_{\rm tot}(R)=A/R^2-B/R$, with
$A>0$, $B>0$. Differentiate:
$\dd E/\dd R=-2A/R^3+B/R^2=0$.
The equilibrium radius is $R=2A/B$:
\begin{equation}
\boxed{R=\left(\frac{9\pi}{4}\right)^{2/3}
\frac{\sh^2 d^{5/3}}{GmM^2N^{1/3}}}.
\end{equation}
At this point $\dd^2E/\dd R^2=2A/R^4>0$,
so it is an energy minimum.
Putting $d=1/2$ and conventional nucleon
masses gives approximately
$\boxed{R=7.6\times10^{25}N^{-1/3}\ {\rm m}}$.
For a solar mass $M_\odot\simeq
1.989\times10^{30}$ kg,
$N\simeq M_\odot/M\simeq1.19\times10^{57}$,
and
\begin{equation}
\boxed{R_{\rm WD}\simeq7.2\times10^3\ {\rm km}}.
\end{equation}
Compute the electron density
$n_e=dN/(4\pi R^3/3)$ and
$E_F=\sh^2(3\pi^2n_e)^{2/3}/(2m)$.
For this model the result is
$\boxed{E_F\sim2.0\times10^5\ {\rm eV}}$,
not dramatically below the electron
rest energy $m_ec^2=5.11\times10^5$ eV.
Thus relativistic corrections are significant
and the nonrelativistic white-dwarf model
is only an approximation.

\pnum{5.36}{Ultrarelativistic gas, critical stellar mass and neutron star}
(a) For $E(k)\simeq\sh ck$,
two spin orientations supply
$V k^2\dd k/\pi^2$ states.
Integrating gives
\begin{equation}
\boxed{E_{\rm ultrarel}
=\frac{V\sh c}{\pi^2}\int_0^{k_F}k^3\dd k
=\frac{V\sh ck_F^4}{4\pi^2}
=\frac34N_e\sh ck_F}.
\end{equation}
(b) With $N_e=dN$, $V=4\pi R^3/3$,
\begin{equation}
E_e^{\rm UR}=
\frac34\left(\frac{9\pi}{4}\right)^{1/3}
\frac{\sh c\,d^{4/3}N^{4/3}}R.
\end{equation}
Add $E_g=-3GM^2N^2/(5R)$. Both terms
scale as $R^{-1}$, so the
nonrelativistic $R^{-2}$ stabilization
disappears. The coefficient of $1/R$
changes sign at
\begin{equation}
\boxed{N_c=\left[
\frac54\left(\frac{9\pi}{4}\right)^{1/3}
\frac{\sh c\,d^{4/3}}{GM^2}
\right]^{3/2}}.
\end{equation}
For $d=1/2$, $M$ equal to a nucleon mass,
the estimate is
$\boxed{N_c\simeq2.04\times10^{57}}$,
equivalent to about
$\boxed{1.7M_\odot}$.
This is the uniform-density Newtonian toy
model's Chandrasekhar-like limit, not a
precision modern white-dwarf mass prediction.
(c) Suppose a solar-mass object is supported
by \emph{neutron} degeneracy rather than
electron degeneracy. Set fermion mass
$m=m_n$, gravitational nucleon mass
$M\simeq m_n$, and one fermion per nucleon
$d=1$ in the nonrelativistic radius formula:
\begin{equation}
\boxed{R_{\rm NS}\sim1.2\times10^4\ {\rm m}
\sim12\ {\rm km}}.
\end{equation}
The density and Fermi energy are
\begin{equation}
n_n=\frac{N}{4\pi R^3/3},\quad
E_{F,n}=\frac{\sh^2}{2m_n}(3\pi^2n_n)^{2/3}
\sim\boxed{6\times10^7\ {\rm eV}}.
\end{equation}
Compare with $m_nc^2\simeq939$ MeV:
the kinetic Fermi scale is roughly
$6\%$ of the rest energy, so an
initial nonrelativistic estimate is not
grossly inconsistent. In real neutron stars,
strong interactions and general relativity
are essential.

\pnum{5.37}{Density of states and one-dimensional van Hove peaks}
A periodic chain has $N$ cells of spacing
$a$, and allowed $q$ spacing
$\Delta q=2\pi/(Na)$. Counting two
momenta $\pm q$ with the same energy,
the DOS per unit length is
\begin{equation}
\boxed{g(E)\equiv\frac{G(E)}{Na}
=\frac1{\pi|\dd E/\dd q|}}.
\end{equation}
(a) For a free particle,
$E=\sh^2q^2/(2m)$, so
\begin{equation}
\boxed{g_{\rm free}(E)
=\frac1{\pi\sh}\sqrt{\frac{m}{2E}}}.
\end{equation}
(b) For the delta lattice let
$k=\sqrt{2mE}/\sh$ and
\begin{equation}
f(k)=\cos(ka)+
\frac{m\alpha}{\sh^2k}\sin(ka),
\qquad \cos(qa)=f(k).
\end{equation}
Differentiate the implicit equation:
$-a\sin(qa)\,\dd q/\dd E
=(\dd f/\dd k)(\dd k/\dd E)$,
where $\dd k/\dd E=m/(\sh^2k)$.
Since $\sin^2(qa)=1-f(k)^2$,
\begin{equation}
\boxed{g_\alpha(E)=
\frac{m}{\pi a\sh^2 k}
\frac{|f'(k)|}{\sqrt{1-f(k)^2}}},
\end{equation}
valid in allowed bands $|f(k)|<1$, with
\begin{equation}
f'(k)=-a\sin(ka)+\frac{m\alpha}{\sh^2}
\left[\frac{a\cos(ka)}k-
\frac{\sin(ka)}{k^2}\right].
\end{equation}
For $|f(k)|>1$ there are no real
Bloch modes and $g_\alpha(E)=0$.
(c) The numerical script compares
$g_0(E)$ and $g_{\alpha=1}(E)$
for $m=\sh=a=1$. Near a band edge,
the denominator $\sqrt{1-f^2}$
vanishes while the derivative is
generally nonzero, producing the
characteristic one-dimensional
van Hove divergence.
\checkit{The DOS formula counts the
two spatial momenta $\pm q$.
If electron spin is also counted,
multiply this spatial DOS by two.}

\pnum{5.38}{Harmonic chain, normal modes and phonon operators}
Consider $N$ equal masses on a ring,
$x_{N+1}=x_1$, with Hamiltonian
\begin{equation}
H=-\frac{\sh^2}{2m}\sum_{j=1}^N
\frac{\partial^2}{\partial x_j^2}
+\frac{m\omega^2}{2}
\sum_{j=1}^N(x_{j+1}-x_j)^2.
\end{equation}
(a) The geometric series satisfies
$\sum_{j=1}^N\ee^{2\pi\ii j n/N}=N$
when $n$ is divisible by $N$ and
zero otherwise. Therefore for
$1\le k,k'\le N-1$,
\begin{equation}
\frac1N\sum_j\ee^{2\pi\ii j(k-k')/N}
=\delta_{kk'},\qquad
\frac1N\sum_j\ee^{2\pi\ii j(k+k')/N}
=\delta_{k',N-k}.
\end{equation}
(b) Define $\omega_k=2\omega\sin(\pi k/N)$,
and the operators
\begin{equation}
a_k^\pm=\frac1{\sqrt N}\sum_{j=1}^N
\ee^{\pm2\pi\ii jk/N}
\left[
\sqrt{\frac{m\omega_k}{2\sh}}\,x_j
\mp\sqrt{\frac{\sh}{2m\omega_k}}\,
\partial_{x_j}\right].
\end{equation}
Using $[\partial_{x_j},x_l]=\delta_{jl}$
and the Fourier orthogonality from (a),
the only surviving terms in
$[a_k^-,a_{k'}^+]$ are the two
coordinate-derivative cross terms.
Their coefficients add to one:
\begin{equation}
\boxed{[a_k^-,a_{k'}^+]=\delta_{kk'},\qquad
[a_k^-,a_{k'}^-]=[a_k^+,a_{k'}^+]=0}.
\end{equation}
(c) Let $R=N^{-1}\sum_jx_j$.
Inverse discrete Fourier transformation,
using (a), gives
\begin{align}
\boxed{x_j}&=R+\frac1{\sqrt N}\sum_{k=1}^{N-1}
\sqrt{\frac{\sh}{2m\omega_k}}
(a_k^-+a_{N-k}^+)\ee^{2\pi\ii jk/N},\\
\boxed{\partial_{x_j}}&=
\frac1N\partial_R+
\frac1{\sqrt N}\sum_{k=1}^{N-1}
\sqrt{\frac{m\omega_k}{2\sh}}
(a_k^--a_{N-k}^+)\ee^{2\pi\ii jk/N}.
\end{align}
(d) Substituting these inverses into
the kinetic and nearest-neighbor
quadratic potential, Fourier orthogonality
eliminates all $k\ne k'$ cross terms.
The result is a free center of mass
plus $N-1$ independent oscillators:
\begin{equation}
\boxed{H=-\frac{\sh^2}{2Nm}\partial_R^2+
\sum_{k=1}^{N-1}\sh\omega_k
\left(a_k^+a_k^-+\frac12\right)}.
\end{equation}
For a center-of-mass plane wave
$\ee^{\ii K R}$ the \emph{energy is positive}:
\begin{equation}
\boxed{E=\frac{\sh^2K^2}{2Nm}+
\sum_{k=1}^{N-1}\sh\omega_k
\left(n_k+\frac12\right)},\quad
n_k=0,1,\dots.
\end{equation}
Each $n_k$ counts phonons in mode $k$.
\checkit{The sign of the free center-of-mass
energy follows from acting
$-\partial_R^2$ on $\ee^{\ii K R}$.
Any negative sign in the printed
energy eigenvalue is a typographic
inconsistency, not a negative kinetic energy.}

\pnum{5.39}{Why periodic and hard-wall free-electron gases agree}
For periodic box lengths $L_x,L_y,L_z$,
the condition
$\psi(\mathbf r+L_i\hat e_i)=\psi(\mathbf r)$
on a plane wave
$\psi=V^{-1/2}\ee^{\ii\mathbf k\cdot\mathbf r}$
forces $\ee^{\ii k_iL_i}=1$.
Therefore
\begin{equation}
\boxed{k_iL_i=2\pi n_i,\qquad
n_i\in\mathbb Z}.
\end{equation}
Each allowed $\mathbf k$ occupies
$\boxed{\Delta k_x\Delta k_y\Delta k_z
=(2\pi)^3/V=8\pi^3/V}$.
This cell is eight times larger
than the $\pi^3/V$ cell of hard-wall
standing-wave states, but periodic
$\mathbf k$ spans the \emph{entire}
sphere rather than only its positive
octant. These factors of eight cancel.
Including the same factor two for spin,
\begin{equation}
N_e=\frac{2V}{(2\pi)^3}
\frac{4\pi k_F^3}{3}
=\frac{Vk_F^3}{3\pi^2},
\end{equation}
so
\begin{equation}
\boxed{k_F=(3\pi^2n)^{1/3},\qquad
E_F=\frac{\sh^2}{2m}(3\pi^2n)^{2/3},\qquad
E_{\rm tot}=\frac35N_eE_F}.
\end{equation}
Bulk thermodynamic quantities therefore
agree in the large-$N$ limit; surface
and finite-size spectra differ.

\section{Numerical verification and figures}
The companion script \texttt{scripts/verify.py} reproduces independent
spectrum, normalization, and band-structure checks using NumPy and SciPy.
\IfFileExists{figures/exchange.pdf}{\begin{figure}[htbp]
\centering\includegraphics[width=.85\linewidth]{figures/exchange.pdf}
\caption{Relative separation density for two oscillator orbitals:
distinguishable, bosonic, and fermionic states (Problem 5.7).}
\end{figure}}{}
\IfFileExists{figures/bands.pdf}{\begin{figure}[htbp]
\centering\includegraphics[width=.85\linewidth]{figures/bands.pdf}
\caption{Allowed Bloch bands of a repulsive periodic delta comb
at $\alpha=m=a=\sh=1$ (Problem 5.28).}
\end{figure}}{}
\IfFileExists{figures/dos.pdf}{\begin{figure}[htbp]
\centering\includegraphics[width=.85\linewidth]{figures/dos.pdf}
\caption{One-dimensional density of states of a free electron
and a repulsive delta-function comb; allowed-band boundaries
generate van Hove peaks (Problem 5.37).}
\end{figure}}{}
\clearpage
\section{Reference and limitations}
D.~J. Griffiths and D.~F. Schroeter,
\emph{Introduction to Quantum Mechanics}, third edition,
Cambridge University Press (2018), Chapter 5;
DOI \href{https://doi.org/10.1017/9781316995433}{10.1017/9781316995433}.
These original independent educational derivations cite textbook
problem numbers only and do not reproduce complete copyrighted prompts.
Simple band, atom, and stellar models omit many-body, relativistic,
surface, or material-specific corrections except where stated.
\end{document}